\section{Homology presentation and decomposition}

Let $K_n$ be the restricted kernel-coordinate matrix.  The kernel-basis
equivalence maps the boundary submodule onto the column span of $K_n$.
Quotient transport therefore gives a linear equivalence
\[
  H_n
  \simeq_{\Z}
  \Z^{\rank C_n-r}/\operatorname{im}K_n.
\]
Lean names this equivalence \code{homologyPresentationEquiv}.  Its codomain
is the \code{presentedModule} of a
\code{FiniteFreePresentation}.

A second strong Smith calculation implements the presentation-matrix step in
the standard Smith computation of homology
\citep[Chapter 3]{kaczynskiMischaikowMrozek2004}:
\[
  U'K_nV'=S'
\]
classifies that cokernel.  If the nonzero normalized entries of $S'$ are
$a_1,\ldots,a_t$, the resulting additive equivalence is
\[
  H_n \simeq
  \left(\prod_{i=1}^{t}\Z/a_i\Z\right)
  \times \Z^{\,(\rank C_n-r)-t}.
\]
The final exponent is \code{bettiNumber}.  The complete factor family keeps
unit entries: $\Z/1\Z$ is a zero summand, and retaining it preserves the
exact certified Smith data.  The reader
\code{homologyTorsionFactors} separately filters units for display.

The runtime projection repeats the same matrix path without evaluating the
noncomputable quotient equivalences.  It reads the outgoing nonzero Smith
prefix, extracts the tail matrix, runs the second strong Smith algorithm, and
returns ranks and factors.  Proofs identify the runtime rank with the
semantic rank, reindex the runtime matrix to $K_n$, and show that its factors
come directly from the second strong result.  Runtime output is therefore a
projection of certified objects rather than an independent calculator.
